\documentclass[letterpaper,10pt]{article}
\usepackage[utf8]{inputenc}
\usepackage[T1]{fontenc}
\usepackage{lmodern}
\usepackage{microtype}
\usepackage[margin=0.45in]{geometry}
\usepackage{hyperref}
\usepackage{titlesec}
\usepackage{enumitem}
\usepackage{xcolor}

\hypersetup{
    colorlinks=true,
    linkcolor=blue,
    urlcolor=blue,
    pdftitle={Saul Cordoba - CV},
    pdfauthor={Saul Cordoba},
    pdfsubject={Software Developer | Backend, DevOps \& AI Engineer},
}

\pagestyle{empty}
\urlstyle{same}

\titleformat{\section}{\large\bfseries}{}{0em}{}[\titlerule]
\titlespacing*{\section}{0pt}{5pt}{3pt}

\setlist[itemize]{leftmargin=1.2em, itemsep=1pt, topsep=1.5pt, parsep=0pt, partopsep=0pt}

\begin{document}

%--- HEADER ---
\begin{center}
    {\Huge \bfseries SAÚL CÓRDOBA} \\[3pt]
    {\Large \textbf{Software Developer | Backend, DevOps \& AI Engineer}} \\[2pt]
    Villavicencio, Colombia (Disponible para Trabajo Remoto / Relocalización) \\
    (+57) 321 961 1182 | \href{mailto:sauleduardo1111@gmail.com}{sauleduardo1111@gmail.com} | \href{https://github.com/Garetetete}{GitHub: Garetetete} | \href{https://www.linkedin.com/in/saul-eduardo-cordoba-romero-13b5b72b6}{LinkedIn} | \textbf{Inglés:} B2
\end{center}

%--- PERFIL PROFESIONAL ---
\section{PERFIL PROFESIONAL}
\textbf{Ingeniero Electrónico} con doble posgrado: \textbf{Especialista en Instrumentación y Control Industrial} y \textbf{Especialista en Ingeniería de Software}. Desarrollador full-stack con foco en \textbf{backend (Python/FastAPI y C\#/.NET), DevOps y sistemas con Inteligencia Artificial generativa (IA / AI)}. \textbf{Diseñé y construí desde cero, en solitario y de extremo a extremo, una plataforma SaaS multi-tenant de análisis documental con Inteligencia Artificial, hoy en producción}: backend, frontend, modelo de datos e infraestructura, con 120+ endpoints REST/SSE y 8 servicios orquestados en Docker. Experiencia adicional en BI (Tableau/Power BI) para el sector petrolero de EE.UU.

%--- HABILIDADES TÉCNICAS ---
\section{HABILIDADES TÉCNICAS}
\begin{itemize}
    \item \textbf{Lenguajes:} Python, C\# (.NET), C, C++, Java, TypeScript/JavaScript, SQL, MATLAB.
    \item \textbf{Backend \& Frontend:} FastAPI, .NET (Clean Architecture), Node/Express, Angular, React, microservicios, API gateways, SSE.
    \item \textbf{DevOps \& Cloud:} Docker y Docker Compose, Linux, SSH/despliegue en servidores, nginx, Keycloak (OAuth2/OIDC), Azure DevOps, GitHub Actions, CI/CD, despliegue continuo.
    \item \textbf{Inteligencia Artificial (IA / AI) y LLMs:} Vertex AI, Gemini, Google ADK (sistemas multiagente), APIs de Claude y OpenAI, RAG con embeddings y búsqueda semántica, salida estructurada, optimización de costos de inferencia.
    \item \textbf{Datos \& BI:} SQL, Oracle, PostgreSQL (RLS, JSONB), SQL Server, MongoDB, Redis, Tableau, Power BI, Looker Studio, Pandas.
    \item \textbf{Metodologías:} SDD (Spec-Driven Development), TDD, Clean Architecture.
\end{itemize}

%--- EXPERIENCIA ---
\section{EXPERIENCIA PROFESIONAL}

\textbf{Desarrollador de Software} \hfill \textbf{2025 -- Presente} \\
\textit{SASOFTCO} \hfill \textit{Villavicencio, Colombia}
\begin{itemize}
    \item \textbf{Creador de una plataforma SaaS de análisis documental con Inteligencia Artificial (IA / AI)} (\textbf{en producción}): la diseñé y desarrollé desde cero de extremo a extremo — FastAPI (Python 3.11), \textbf{120+ endpoints REST/SSE}, PostgreSQL 15 \textbf{multi-tenant con seguridad a nivel de fila (RLS)} y 8 servicios en Docker Compose, incluido un microservicio C\#/.NET.
    \item \textbf{Inteligencia Artificial generativa y sistemas multiagente (AI/LLM)}: orquestación de agentes especializados sobre \textbf{Vertex AI y Gemini}, con salida estructurada en JSON y capa de abstracción multiproveedor (Gemini, Claude, OpenAI).
    \item \textbf{Reducción superior al 50\% en costos de inferencia de IA}, mediante control de consumo, caché y tolerancia a fallos en entornos de inferencia distribuidos.
    \item \textbf{Procesamiento asíncrono y en tiempo real}: streaming de progreso vía SSE, procesamiento por lotes con control de concurrencia y reintentos con backoff exponencial; ingesta automatizada desde APIs de datos abiertos gubernamentales.
    \item Desarrollo backend y frontend en \textbf{microservicios .NET con Angular y React} para sector gobierno y asegurador, con BPMN en Camunda y n8n; monolito Java/Spring con integración SAP; pipelines CI/CD en Azure DevOps.
    \item \textbf{Tecnologías usadas:} Inteligencia Artificial (IA/AI), LLM, Vertex AI, Gemini, Google ADK, Python, FastAPI, C\#/.NET, Angular, React, TypeScript, PostgreSQL, Docker, Keycloak, Azure DevOps, GitHub Actions.
\end{itemize}

\textbf{Desarrollador / Analista de Datos} \hfill \textbf{2024 -- 2025} \\
\textit{Btilink SAS} \hfill \textit{Villavicencio, Colombia}
\begin{itemize}
    \item \textbf{Analista de datos para Aera Energy LLC} (compañía petrolera de EE.UU.): consultas \textbf{SQL sobre bases de datos Oracle} para extraer y consolidar información de la operación petrolera.
    \item Dashboards y reportes en \textbf{Tableau y Power BI} (volúmenes de producción, ingresos y salidas) utilizados por varios equipos para el \textbf{monitoreo continuo de la operación}.
\end{itemize}

%--- EDUCACIÓN ---
\section{EDUCACIÓN}
\textbf{Especialización en Ingeniería de Software} \hfill \textbf{2025 -- 2026} \\
Universidad de los Llanos \\
\textbf{Especialización en Instrumentación y Control Industrial} \hfill \textbf{2024 -- 2025} \\
Universidad de los Llanos --- \textit{Modalidad de grado: Diplomado en Sistemas IoT y PLC} \\
\textbf{Ingeniería Electrónica} \hfill \textbf{2018 -- 2024} \\
Universidad de los Llanos

%--- CERTIFICACIONES (descomentar y completar cuando tengas los datos) ---
% \section{CERTIFICACIONES Y DISTINCIONES}
% \begin{itemize}
%     \item \textbf{Nombre de la certificación} -- Entidad emisora (Año).
% \end{itemize}

\end{document}
